% --- Archon physics formalization source begin ---
% archon:physics
% archon:covers PhyXMiniProblems/problem_phyx_mini_0541.lean
% archon:source-report reports/phyx_mini/problem_phyx_mini_0541.source.json

\chapter{Physics problem phyx\_mini\_0541}
\label{ch:PhyXMiniProblems_problem_phyx_mini_0541}

\paragraph{Problem source.}
\#\# Physical scenario

As shown in figure

\#\# Question

Find the probabilities of the measurements shown in figure, assuming that the first Stern-Gerlach analyzer is aligned along the direction \textbackslash{}(\textbackslash{}hat\{\textbackslash{}mathbf\{n\}\}\textbackslash{}) defined by the angles \textbackslash{}(\textbackslash{}theta = 2\textbackslash{}pi/3\textbackslash{}) and \textbackslash{}(\textbackslash{}phi = \textbackslash{}pi/4\textbackslash{}).

\#\# Answer choices

- A: A: \textbackslash{}(\textbackslash{}mathcal\{P\}\_\{+x\} \textbackslash{}approx 0.688 \textbackslash{}mathcal\{P\}\_\{-x\} \textbackslash{}approx 0.312\textbackslash{})

- B: B: \textbackslash{}(\textbackslash{}mathcal\{P\}\_\{+x\} \textbackslash{}approx 0.883 \textbackslash{}mathcal\{P\}\_\{-x\} \textbackslash{}approx 0.117\textbackslash{})

- C: C: \textbackslash{}(\textbackslash{}mathcal\{P\}\_\{+x\} \textbackslash{}approx 0.765 \textbackslash{}mathcal\{P\}\_\{-x\} \textbackslash{}approx 0.235\textbackslash{})

- D: D: \textbackslash{}(\textbackslash{}mathcal\{P\}\_\{+x\} \textbackslash{}approx 0.806 \textbackslash{}mathcal\{P\}\_\{-x\} \textbackslash{}approx 0.194\textbackslash{})

\#\# Figure caption (auxiliary; use the image as primary evidence)

The image illustrates a relativistic system involving several objects and reference frames. Here's a detailed description:

- **Space Station**: Located at the lower left, labeled as "Space station," part of "System K at rest." It appears to be in a stationary state in this reference frame.
  
- **Coordinate Axes**: There are two sets of axes:
  - The **K** system, labeled with \textbackslash{}(x\textbackslash{}) and \textbackslash{}(y\textbackslash{}) axes.
  - The **K'** system, labeled with \textbackslash{}(x'\textbackslash{}) and \textbackslash{}(y'\textbackslash{}) axes, moving with velocity \textbackslash{}(\textbackslash{}vec\{v\}\textbackslash{}).

- **Spaceship**: Positioned along the x-axis, labeled as "Spaceship," moving to the right with velocity \textbackslash{}(\textbackslash{}vec\{v\}\textbackslash{}) in the \textbackslash{}(K'\textbackslash{}) frame. The spaceship is depicted as emitting a trail, indicating motion.

- **Proton**: Located to the right of the spaceship, labeled as "Proton," with two velocities indicated:
  - \textbackslash{}(u\textbackslash{}): In system K.
  - \textbackslash{}(u'\textbackslash{}): In system K'.

- **Asteroid**: Further to the right along the x-axis, labeled "Asteroid," depicted as a blue, irregular shape.

- **Velocity Vectors**: The vectors are denoted by blue arrows:
  - \textbackslash{}(\textbackslash{}vec\{v\}\textbackslash{}): indicating the velocity of the spaceship and the K' system.
  - \textbackslash{}(u\textbackslash{}) and \textbackslash{}(u'\textbackslash{}): representing proton velocities in different frames.

The image seems to be illustrating the concept of relative motion and velocities in special relativity, highlighting differences in observation from various frames.

\#\# Dataset metadata

Relativity; Physical Model Grounding Reasoning, Spatial Relation Reasoning

\paragraph{Recorded answer/context.}
D: D: \textbackslash{}(\textbackslash{}mathcal\{P\}\_\{+x\} \textbackslash{}approx 0.806 \textbackslash{}mathcal\{P\}\_\{-x\} \textbackslash{}approx 0.194\textbackslash{})

\paragraph{Figure/image path.}
/root/proposal\_for\_physic/science-mango/phyx\_mini\_run/phyx\_data/test\_image/541.png

\paragraph{Formalization target.}
create a compiling Lean file with sorry bodies at `PhyXMiniProblems/problem\_phyx\_mini\_0541.lean`. The Lean declarations must preserve the physical quantities, dimensions or dimensional roles, figure labels, governing-law hypotheses, and final relation expressed by this problem.
Use Mathlib/Physlib names found through LeanExplore where available. If a physics API is missing, introduce faithful local abstractions rather than scalar placeholder aliases.

\begin{theorem}[Physics formalization target]
\label{thm:physics:phyx_mini_0541:target}
\lean{PhyXMiniProblems.ProblemPhyXMini0541.problem_phyx_mini_0541}
\uses{def:physics:phyx-mini-0541:phyxminiproblems-problemphyxmini0541-analyzeranglereadouts, def:physics:phyx-mini-0541:phyxminiproblems-problemphyxmini0541-sequentialsterngerlachexperiment, def:physics:phyx-mini-0541:phyxminiproblems-problemphyxmini0541-satisfiesspinhalfbornrule, def:physics:phyx-mini-0541:phyxminiproblems-problemphyxmini0541-satisfiessequentialanalyzersemantics, def:physics:phyx-mini-0541:phyxminiproblems-problemphyxmini0541-matchessterngerlachquestion, def:physics:phyx-mini-0541:phyxminiproblems-problemphyxmini0541-suppliedrelativityraster, def:physics:phyx-mini-0541:phyxminiproblems-problemphyxmini0541-matchessuppliedrelativityraster, def:physics:phyx-mini-0541:phyxminiproblems-problemphyxmini0541-probabilityapproximatelyequal, def:physics:phyx-mini-0541:phyxminiproblems-problemphyxmini0541-isuniquematchinganswerchoice}
The assigned autoformalize agent should translate this physics problem into Lean declarations in the covered file, with theorem and lemma proof bodies written as `by sorry`.
\end{theorem}
\begin{proof}
This is an autoformalization task, not a proof task. Produce statements that can later be proved by the physics prover without weakening the physical model.
\end{proof}

% --- Archon declaration topology begin ---
\section{Declaration topology}
\begin{definition}[Spin Direction]
  \label{def:physics:phyx-mini-0541:phyxminiproblems-problemphyxmini0541-spindirection}
  \lean{PhyXMiniProblems.ProblemPhyXMini0541.SpinDirection}
  A physical spin/analyzer direction, grounded as the unit metric sphere in three-dimensional Euclidean space
\end{definition}
\begin{proof}
  This is the declared model definition of Spin Direction.
\end{proof}
\begin{definition}[Spin Outcome]
  \label{def:physics:phyx-mini-0541:phyxminiproblems-problemphyxmini0541-spinoutcome}
  \lean{PhyXMiniProblems.ProblemPhyXMini0541.SpinOutcome}
  The two output ports of a spin-one-half Stern--Gerlach analyzer
\end{definition}
\begin{proof}
  This is the declared model definition of Spin Outcome.
\end{proof}
\begin{definition}[outcome Sign]
  \label{def:physics:phyx-mini-0541:phyxminiproblems-problemphyxmini0541-outcomesign}
  \lean{PhyXMiniProblems.ProblemPhyXMini0541.outcomeSign}
  \uses{def:physics:phyx-mini-0541:phyxminiproblems-problemphyxmini0541-spinoutcome}
  The eigenvalue sign associated with an analyzer output
\end{definition}
\begin{proof}
  This is the declared model definition of outcome Sign.
\end{proof}
\begin{definition}[spin Direction Dot]
  \label{def:physics:phyx-mini-0541:phyxminiproblems-problemphyxmini0541-spindirectiondot}
  \lean{PhyXMiniProblems.ProblemPhyXMini0541.spinDirectionDot}
  \uses{def:physics:phyx-mini-0541:phyxminiproblems-problemphyxmini0541-spindirection}
  Euclidean scalar product of two physical analyzer directions
\end{definition}
\begin{proof}
  This is the declared model definition of spin Direction Dot.
\end{proof}
\begin{definition}[spherical Spin Direction]
  \label{def:physics:phyx-mini-0541:phyxminiproblems-problemphyxmini0541-sphericalspindirection}
  \lean{PhyXMiniProblems.ProblemPhyXMini0541.sphericalSpinDirection}
  \uses{def:physics:phyx-mini-0541:phyxminiproblems-problemphyxmini0541-spindirection}
  The unit direction with polar angle thetaRadians and azimuthal angle phiRadians, in the usual (x,y,z) spherical-coordinate convention
\end{definition}
\begin{proof}
  This is the declared model definition of spherical Spin Direction.
\end{proof}
\begin{definition}[positive XDirection]
  \label{def:physics:phyx-mini-0541:phyxminiproblems-problemphyxmini0541-positivexdirection}
  \lean{PhyXMiniProblems.ProblemPhyXMini0541.positiveXDirection}
  \uses{def:physics:phyx-mini-0541:phyxminiproblems-problemphyxmini0541-spindirection}
  The positive Cartesian x direction used by the second analyzer
\end{definition}
\begin{proof}
  This is the declared model definition of positive XDirection.
\end{proof}
\begin{definition}[Analyzer Angle Readouts]
  \label{def:physics:phyx-mini-0541:phyxminiproblems-problemphyxmini0541-analyzeranglereadouts}
  \lean{PhyXMiniProblems.ProblemPhyXMini0541.AnalyzerAngleReadouts}
  Scalar angle readouts printed in the question, expressed in radians
\end{definition}
\begin{proof}
  This is the declared model definition of Analyzer Angle Readouts.
\end{proof}
\begin{definition}[Spin Half Transition Model]
  \label{def:physics:phyx-mini-0541:phyxminiproblems-problemphyxmini0541-spinhalftransitionmodel}
  \lean{PhyXMiniProblems.ProblemPhyXMini0541.SpinHalfTransitionModel}
  \uses{def:physics:phyx-mini-0541:phyxminiproblems-problemphyxmini0541-spindirection, def:physics:phyx-mini-0541:phyxminiproblems-problemphyxmini0541-spinoutcome}
  Transition probabilities between analyzer outputs. The arguments are, in order, the preparation axis, selected preparation outcome, measurement axis, and measured outcome. The returned real number is a dimensionless probability readout
\end{definition}
\begin{proof}
  This is the declared model definition of Spin Half Transition Model.
\end{proof}
\begin{definition}[Sequential Stern Gerlach Experiment]
  \label{def:physics:phyx-mini-0541:phyxminiproblems-problemphyxmini0541-sequentialsterngerlachexperiment}
  \lean{PhyXMiniProblems.ProblemPhyXMini0541.SequentialSternGerlachExperiment}
  \uses{def:physics:phyx-mini-0541:phyxminiproblems-problemphyxmini0541-spindirection, def:physics:phyx-mini-0541:phyxminiproblems-problemphyxmini0541-spinoutcome, def:physics:phyx-mini-0541:phyxminiproblems-problemphyxmini0541-spinhalftransitionmodel}
  A sequential two-analyzer experiment. The selected output of the first analyzer prepares the beam incident on the second analyzer. secondOutcomeProbabilityReadout is kept distinct from the abstract transition model until the apparatus-semantics premise relates them
\end{definition}
\begin{proof}
  This is the declared model definition of Sequential Stern Gerlach Experiment.
\end{proof}
\begin{definition}[Satisfies Spin Half Born Rule]
  \label{def:physics:phyx-mini-0541:phyxminiproblems-problemphyxmini0541-satisfiesspinhalfbornrule}
  \lean{PhyXMiniProblems.ProblemPhyXMini0541.SatisfiesSpinHalfBornRule}
  \uses{def:physics:phyx-mini-0541:phyxminiproblems-problemphyxmini0541-outcomesign, def:physics:phyx-mini-0541:phyxminiproblems-problemphyxmini0541-spindirectiondot, def:physics:phyx-mini-0541:phyxminiproblems-problemphyxmini0541-spinhalftransitionmodel}
  General spin-one-half Born transition law for arbitrary preparation and measurement axes and for both output signs. This law does not mention the angles or numerical probabilities of the current question
\end{definition}
\begin{proof}
  This is the declared model definition of Satisfies Spin Half Born Rule.
\end{proof}
\begin{definition}[Satisfies Sequential Analyzer Semantics]
  \label{def:physics:phyx-mini-0541:phyxminiproblems-problemphyxmini0541-satisfiessequentialanalyzersemantics}
  \lean{PhyXMiniProblems.ProblemPhyXMini0541.SatisfiesSequentialAnalyzerSemantics}
  \uses{def:physics:phyx-mini-0541:phyxminiproblems-problemphyxmini0541-sequentialsterngerlachexperiment}
  The apparatus readout is the transition probability from the selected first-analyzer branch to the indicated second-analyzer branch
\end{definition}
\begin{proof}
  This is the declared model definition of Satisfies Sequential Analyzer Semantics.
\end{proof}
\begin{definition}[Matches Stern Gerlach Question]
  \label{def:physics:phyx-mini-0541:phyxminiproblems-problemphyxmini0541-matchessterngerlachquestion}
  \lean{PhyXMiniProblems.ProblemPhyXMini0541.MatchesSternGerlachQuestion}
  \uses{def:physics:phyx-mini-0541:phyxminiproblems-problemphyxmini0541-sphericalspindirection, def:physics:phyx-mini-0541:phyxminiproblems-problemphyxmini0541-positivexdirection, def:physics:phyx-mini-0541:phyxminiproblems-problemphyxmini0541-analyzeranglereadouts, def:physics:phyx-mini-0541:phyxminiproblems-problemphyxmini0541-sequentialsterngerlachexperiment}
  The prose readouts and analyzer arrangement in the quantum question. The first analyzer selects the positive branch along the specified spherical axis; the second analyzer measures along positive x. No target probability appears here
\end{definition}
\begin{proof}
  This is the declared model definition of Matches Stern Gerlach Question.
\end{proof}
\begin{definition}[Raster Frame]
  \label{def:physics:phyx-mini-0541:phyxminiproblems-problemphyxmini0541-rasterframe}
  \lean{PhyXMiniProblems.ProblemPhyXMini0541.RasterFrame}
  Reference-frame labels visible in the supplied raster
\end{definition}
\begin{proof}
  This is the declared model definition of Raster Frame.
\end{proof}
\begin{definition}[Raster Object]
  \label{def:physics:phyx-mini-0541:phyxminiproblems-problemphyxmini0541-rasterobject}
  \lean{PhyXMiniProblems.ProblemPhyXMini0541.RasterObject}
  Object labels visible from left to right in the supplied raster
\end{definition}
\begin{proof}
  This is the declared model definition of Raster Object.
\end{proof}
\begin{definition}[Raster Axis]
  \label{def:physics:phyx-mini-0541:phyxminiproblems-problemphyxmini0541-rasteraxis}
  \lean{PhyXMiniProblems.ProblemPhyXMini0541.RasterAxis}
  Coordinate-axis labels visible in the supplied raster
\end{definition}
\begin{proof}
  This is the declared model definition of Raster Axis.
\end{proof}
\begin{definition}[Raster Velocity Symbol]
  \label{def:physics:phyx-mini-0541:phyxminiproblems-problemphyxmini0541-rastervelocitysymbol}
  \lean{PhyXMiniProblems.ProblemPhyXMini0541.RasterVelocitySymbol}
  Velocity symbols visible in the supplied raster
\end{definition}
\begin{proof}
  This is the declared model definition of Raster Velocity Symbol.
\end{proof}
\begin{definition}[Supplied Relativity Raster]
  \label{def:physics:phyx-mini-0541:phyxminiproblems-problemphyxmini0541-suppliedrelativityraster}
  \lean{PhyXMiniProblems.ProblemPhyXMini0541.SuppliedRelativityRaster}
  \uses{def:physics:phyx-mini-0541:phyxminiproblems-problemphyxmini0541-rasterframe, def:physics:phyx-mini-0541:phyxminiproblems-problemphyxmini0541-rasterobject, def:physics:phyx-mini-0541:phyxminiproblems-problemphyxmini0541-rasteraxis, def:physics:phyx-mini-0541:phyxminiproblems-problemphyxmini0541-rastervelocitysymbol}
  Calibrated readouts from the supplied image artifact. Horizontal coordinates are dimensionless image-coordinate readouts used only for ordering. Velocity arrows are two-dimensional coordinate readouts in an unspecified common speed unit, since the raster prints no numerical scale
\end{definition}
\begin{proof}
  This is the declared model definition of Supplied Relativity Raster.
\end{proof}
\begin{definition}[Points Along Positive Horizontal Axis]
  \label{def:physics:phyx-mini-0541:phyxminiproblems-problemphyxmini0541-pointsalongpositivehorizontalaxis}
  \lean{PhyXMiniProblems.ProblemPhyXMini0541.PointsAlongPositiveHorizontalAxis}
  A two-dimensional arrow points strictly in the positive horizontal direction
\end{definition}
\begin{proof}
  This is the declared model definition of Points Along Positive Horizontal Axis.
\end{proof}
\begin{definition}[Matches Supplied Relativity Raster]
  \label{def:physics:phyx-mini-0541:phyxminiproblems-problemphyxmini0541-matchessuppliedrelativityraster}
  \lean{PhyXMiniProblems.ProblemPhyXMini0541.MatchesSuppliedRelativityRaster}
  \uses{def:physics:phyx-mini-0541:phyxminiproblems-problemphyxmini0541-suppliedrelativityraster, def:physics:phyx-mini-0541:phyxminiproblems-problemphyxmini0541-pointsalongpositivehorizontalaxis}
  Literal labels and spatial relations read from the supplied raster. The last field explicitly records the source mismatch: the image is a relativity illustration rather than the analyzer diagram described by the question
\end{definition}
\begin{proof}
  This is the declared model definition of Matches Supplied Relativity Raster.
\end{proof}
\begin{definition}[Answer Choice]
  \label{def:physics:phyx-mini-0541:phyxminiproblems-problemphyxmini0541-answerchoice}
  \lean{PhyXMiniProblems.ProblemPhyXMini0541.AnswerChoice}
  Labels of the four answer choices printed with the question
\end{definition}
\begin{proof}
  This is the declared model definition of Answer Choice.
\end{proof}
\begin{definition}[displayed Probability Pair]
  \label{def:physics:phyx-mini-0541:phyxminiproblems-problemphyxmini0541-displayedprobabilitypair}
  \lean{PhyXMiniProblems.ProblemPhyXMini0541.displayedProbabilityPair}
  \uses{def:physics:phyx-mini-0541:phyxminiproblems-problemphyxmini0541-answerchoice}
  The displayed (P\_\{+x\}, P\_\{-x\}) pair for each answer choice
\end{definition}
\begin{proof}
  This is the declared model definition of displayed Probability Pair.
\end{proof}
\begin{definition}[recorded Dataset Answer]
  \label{def:physics:phyx-mini-0541:phyxminiproblems-problemphyxmini0541-recordeddatasetanswer}
  \lean{PhyXMiniProblems.ProblemPhyXMini0541.recordedDatasetAnswer}
  \uses{def:physics:phyx-mini-0541:phyxminiproblems-problemphyxmini0541-answerchoice}
  The answer recorded in the dataset, retained only as source metadata
\end{definition}
\begin{proof}
  This is the declared model definition of recorded Dataset Answer.
\end{proof}
\begin{definition}[Probability Approximately Equal]
  \label{def:physics:phyx-mini-0541:phyxminiproblems-problemphyxmini0541-probabilityapproximatelyequal}
  \lean{PhyXMiniProblems.ProblemPhyXMini0541.ProbabilityApproximatelyEqual}
  Absolute-error comparison for dimensionless probability readouts
\end{definition}
\begin{proof}
  This is the declared model definition of Probability Approximately Equal.
\end{proof}
\begin{definition}[Matches Displayed Choice Within]
  \label{def:physics:phyx-mini-0541:phyxminiproblems-problemphyxmini0541-matchesdisplayedchoicewithin}
  \lean{PhyXMiniProblems.ProblemPhyXMini0541.MatchesDisplayedChoiceWithin}
  \uses{def:physics:phyx-mini-0541:phyxminiproblems-problemphyxmini0541-sequentialsterngerlachexperiment, def:physics:phyx-mini-0541:phyxminiproblems-problemphyxmini0541-answerchoice, def:physics:phyx-mini-0541:phyxminiproblems-problemphyxmini0541-displayedprobabilitypair, def:physics:phyx-mini-0541:phyxminiproblems-problemphyxmini0541-probabilityapproximatelyequal}
  A displayed pair agrees with both modeled branch probabilities to the chosen absolute tolerance
\end{definition}
\begin{proof}
  This is the declared model definition of Matches Displayed Choice Within.
\end{proof}
\begin{definition}[Is Unique Matching Answer Choice]
  \label{def:physics:phyx-mini-0541:phyxminiproblems-problemphyxmini0541-isuniquematchinganswerchoice}
  \lean{PhyXMiniProblems.ProblemPhyXMini0541.IsUniqueMatchingAnswerChoice}
  \uses{def:physics:phyx-mini-0541:phyxminiproblems-problemphyxmini0541-sequentialsterngerlachexperiment, def:physics:phyx-mini-0541:phyxminiproblems-problemphyxmini0541-answerchoice, def:physics:phyx-mini-0541:phyxminiproblems-problemphyxmini0541-matchesdisplayedchoicewithin}
  A choice is the unique displayed pair agreeing with the model
\end{definition}
\begin{proof}
  This is the declared model definition of Is Unique Matching Answer Choice.
\end{proof}
\begin{lemma}[specified Direction dot positive X]
  \label{lem:physics:phyx-mini-0541:phyxminiproblems-problemphyxmini0541-specifieddirection-dot-positivex}
  \lean{PhyXMiniProblems.ProblemPhyXMini0541.specifiedDirection_dot_positiveX}
  \uses{def:physics:phyx-mini-0541:phyxminiproblems-problemphyxmini0541-spindirectiondot, def:physics:phyx-mini-0541:phyxminiproblems-problemphyxmini0541-sphericalspindirection, def:physics:phyx-mini-0541:phyxminiproblems-problemphyxmini0541-positivexdirection}
  The stated angles give n · x = sqrt 6 / 4 in the adopted spherical coordinate convention
\end{lemma}
\begin{proof}
  The conclusion follows from the governing relations and figure readouts named in the statement of specified Direction dot positive X.
\end{proof}
% --- Archon declaration topology end ---
% --- Archon physics formalization source end ---
